\fancyhead[L]{\small SEYFFER ÜBER DIE GEOGRAFISCHE LÄNGE.}
\hypertarget{native-538Seyffer}{}\label{native-538Seyffer}
\begin{center}\small Göttingische gelehrte Anzeigen. Stück 78. S. 773..775. 18. Mai 1811.\\[3pt]\rule{.72\linewidth}{.3pt}\end{center}

\emph{Über die geografische Länge der königlichen Sternwarte in München, erstmals bestimmt aus siebenunddreißig beobachteten und rechnerisch ausgewerteten Sonnenfinsternissen, von} \textsc{Carl Felix Seyffer}. 1810. 104 Seiten im Quartformat. München.

Herr \textsc{Seyffer} leitet in dieser Abhandlung, die aus den Denkschriften der königlichen Bayerischen Akademie der Wissenschaften gesondert abgedruckt wurde, die Länge von München her. Der Titel lässt auf 37 Sonnenfinsternisse schließen. Tatsächlich vergleicht Seyffer die in München gemachte Beobachtung einer einzigen Sonnenfinsternis (vom 16. Juni 1806) mit 36 Beobachtungen an anderen Orten. Verglichen wurden die Orte \emph{Rom} (zwei Beobachtungen), \emph{Padua, Mailand, Madrid, Aranjues, Pampelona, Rinderhook, Fort Orange} (beide in den Vereinigten Staaten von Nordamerika), \emph{Amsterdam, Utrecht,} \SourceLeaf{539}\emph{Zürich, Ochsenhausen, Leipzig, Breslau, Ofen, Cracau, Erlau} (in Ungarn), \emph{Schweidnitz, Hamburg, Luck} (in Polen), \emph{Bourg en Bresse}, die Insel \emph{Leon} (bei Cadiz), \emph{Montauban, Toulouse, Paris, Prag, Lilienthal, Reikevich} (auf Island), \emph{Göttingen, Neapel, Brünn. Berlin, Regensburg, Cremsmünster.} Soweit es allein um die Bestimmung der Länge von München ging, wäre es freilich nicht nötig gewesen, so viele Orte in die Rechnung einzubeziehen. Als Beitrag zur Bestimmung der Längen dieser Orte selbst ist diese Arbeit jedoch wertvoll. Herr \textsc{Seyffer} hat alle Rechnungen in sämtlichen Einzelheiten abdrucken lassen.

Nachdem Herr \textsc{Seyffer} anhand der vollständigen Beobachtungen die Korrekturen der Mondbreite, der Horizontalparallaxe und der Summe der Halbmesser untersucht hat, findet er für die Länge von München als Mittelwert aus 15 Vergleichen $37^{\mathrm m}5^{\mathrm s}6$ östlich von Paris. Im neuesten Band der Connoissance des tems wird München unter den Orten aufgeführt, deren Längen Herr \textsc{Burckhardt} erneut untersucht hat. Dort ist die Länge mit $37^{\mathrm m}0^{\mathrm s}$ angesetzt.
